\documentclass[final]{article}


\usepackage[main]{APS360}
\usepackage[utf8]{inputenc} % allow utf-8 input
\usepackage[T1]{fontenc}    % use 8-bit T1 fonts
\usepackage{hyperref}       % hyperlinks
\usepackage{url}            % simple URL typesetting
\usepackage{booktabs}       % professional-quality tables
\usepackage{amsfonts}       % blackboard math symbols
\usepackage{nicefrac}       % compact symbols for 1/2, etc.
\usepackage{microtype}      % microtypography
\usepackage{xcolor}         % colors
 
\title{Herb Image Classification}

\author{%
  Samantha Hossain \\
  1008967138, samantha.hossain@utoronto.ca\\
  Link to Google Colab\\
}

\begin{document}

\maketitle

\vspace{-0.5in}

\section*{Introduction}
As I started getting into gardening, I realized that I was not able to diffrentiate between different types of herbs. 
I wanted to create a model that could classify different types of herbs based on images. This would be useful for those who just started cooking with herbs, and do not know the difference between all of them to identify herbs quickly and accurately.

\section*{Illustration}
This is how you can reference Figure \ref{fig:sample}.

\begin{figure}[!h]
  \centering
  \fbox{\rule[-.5cm]{0cm}{4cm} \rule[-.5cm]{4cm}{0cm}}
  \caption{Sample figure caption.}
  \label{fig:sample}
\end{figure}

\section*{Background \& Related Work}
Do some background research and cite relevant papers. For example, you can cite a paper like this \cite{DBLP:journals/corr/abs-1802-05365}.

\section*{Data Processing}

\section*{Architecture}

\section*{Baseline Model}

\section*{Ethical Considerations}

\newpage
\section*{References}


\end{document}